\documentclass[10pt,a4paper]{amsart}
\usepackage{amsmath,amssymb,amsthm,graphicx}
\usepackage{verbatim}
\newenvironment{reference}{\comment}{\endcomment}
%\newenvironment{reference}{}{}
\newenvironment{slogan}{\comment}{\endcomment}
\newenvironment{history}{\comment}{\endcomment}

% For commutative diagrams we use Xy-pic
\usepackage[all]{xy}

% We use 2cell for 2-commutative diagrams.
\xyoption{2cell}
% Theorem environments.
%
\theoremstyle{plain}
\newtheorem{theorem}[subsection]{Theorem}
\newtheorem{proposition}[subsection]{Proposition}
\newtheorem{lemma}[subsection]{Lemma}

\theoremstyle{definition}
\newtheorem{definition}[subsection]{Definition}
\newtheorem{example}[subsection]{Example}
\newtheorem{exercise}[subsection]{Exercise}
\newtheorem{situation}[subsection]{Situation}

\theoremstyle{remark}
\newtheorem{remark}[subsection]{Remark}
\newtheorem{remarks}[subsection]{Remarks}

\numberwithin{equation}{subsection}

% Macros
%
\def\lim{\mathop{\mathrm{lim}}\nolimits}
\def\colim{\mathop{\mathrm{colim}}\nolimits}
\def\Spec{\mathop{\mathrm{Spec}}}
\def\Hom{\mathop{\mathrm{Hom}}\nolimits}
\def\Ext{\mathop{\mathrm{Ext}}\nolimits}
\def\SheafHom{\mathop{\mathcal{H}\!\mathit{om}}\nolimits}
\def\SheafExt{\mathop{\mathcal{E}\!\mathit{xt}}\nolimits}
\def\Sch{\mathit{Sch}}
\def\Mor{\mathop{\mathrm{Mor}}\nolimits}
\def\Ob{\mathop{\mathrm{Ob}}\nolimits}
\def\Sh{\mathop{\mathit{Sh}}\nolimits}
\def\NL{\mathop{N\!L}\nolimits}
\def\CH{\mathop{\mathrm{CH}}\nolimits}
\def\proetale{{pro\text{-}\acute{e}tale}}
\def\etale{{\acute{e}tale}}
\def\QCoh{\mathit{QCoh}}
\def\Ker{\mathop{\mathrm{Ker}}}
\def\Im{\mathop{\mathrm{Im}}}
\def\Coker{\mathop{\mathrm{Coker}}}
\def\Coim{\mathop{\mathrm{Coim}}}

% Boxtimes
%
\DeclareMathSymbol{\boxtimes}{\mathbin}{AMSa}{"02}

%
% Macros for moduli stacks/spaces
%
\def\QCohstack{\mathcal{QC}\!\mathit{oh}}
\def\Cohstack{\mathcal{C}\!\mathit{oh}}
\def\Spacesstack{\mathcal{S}\!\mathit{paces}}
\def\Quotfunctor{\mathrm{Quot}}
\def\Hilbfunctor{\mathrm{Hilb}}
\def\Curvesstack{\mathcal{C}\!\mathit{urves}}
\def\Polarizedstack{\mathcal{P}\!\mathit{olarized}}
\def\Complexesstack{\mathcal{C}\!\mathit{omplexes}}
% \Pic is the operator that assigns to X its picard group, usage \Pic(X)
% \Picardstack_{X/B} denotes the Picard stack of X over B
% \Picardfunctor_{X/B} denotes the Picard functor of X over B
\def\Pic{\mathop{\mathrm{Pic}}\nolimits}
\def\Picardstack{\mathcal{P}\!\mathit{ic}}
\def\Picardfunctor{\mathrm{Pic}}
\def\Deformationcategory{\mathcal{D}\!\mathit{ef}}
\usepackage{cleveref}
\crefname{originallemma}{original lemma}{original lemmas}
\Crefname{originallemma}{Original lemma}{Original lemmas}
\makeatletter
\newlabel{lemma-quasi-compact-existence-universally-closed}{{03KA}{source}{}{sourcepointer-1}{}}
\newlabel{lemma-quasi-compact-existence-universally-closed@cref}{{[originallemma][0][]03KA}{[1][0][]source}{}{}{}}
\makeatother
\crefname{originallemma}{original lemma}{original lemmas}
\Crefname{originallemma}{Original lemma}{Original lemmas}
\makeatletter
\newlabel{more-morphisms-lemma-characterize-finite}{{02LS}{source}{}{sourcepointer-2}{}}
\newlabel{more-morphisms-lemma-characterize-finite@cref}{{[originallemma][0][]02LS}{[1][0][]source}{}{}{}}
\makeatother
\crefname{originallemma}{original lemma}{original lemmas}
\Crefname{originallemma}{Original lemma}{Original lemmas}
\makeatletter
\newlabel{morphisms-lemma-finite-monomorphism-closed}{{03BB}{source}{}{sourcepointer-3}{}}
\newlabel{morphisms-lemma-finite-monomorphism-closed@cref}{{[originallemma][0][]03BB}{[1][0][]source}{}{}{}}
\makeatother
\crefname{originallemma}{original lemma}{original lemmas}
\Crefname{originallemma}{Original lemma}{Original lemmas}
\makeatletter
\newlabel{spaces-lemma-representable-transformations-property-implication}{{02YO}{source}{}{sourcepointer-4}{}}
\newlabel{spaces-lemma-representable-transformations-property-implication@cref}{{[originallemma][0][]02YO}{[1][0][]source}{}{}{}}
\makeatother
\title{The valuative criterion for separatedness of morphisms of algebraic spaces}
\author{The Stacks Project Authors}
\date{stacks-project:spaces-morphisms.tex:local-snapshot-20261012; GNU FDL 1.2 or later}
\begin{document}
\maketitle
\setcounter{subsection}{1}
\noindent\textit{Editorial note.} All delivered mathematical text is the author's own, copied from the pinned Stacks Project chapter file spaces-morphisms.tex; no statement or proof step was rewritten, reordered or removed.\par\medskip
\noindent\textit{Editorial note.} The statement carries an optional name argument ('Valuative criterion separatedness'); the optional argument is part of the counted statement span and is preserved verbatim.\par\medskip
\noindent\textit{Editorial note.} The four references of the counted spans are recorded as non-load-bearing navigation pointers; the argument that consumes them stays inside the counted proof span.\par\medskip
\noindent\textit{Editorial note.} Preamble: the three native excerpt spans of sources/stacks-spaces-morphisms/preamble.tex (lines 8-13, 14-18 and 154-217). The chapter file itself carries no preamble text.\par\medskip
\noindent\textit{Allowed standard prerequisites (editorial).}\begin{itemize}
\item Quasi-separatedness of f is quasi-compactness of the diagonal. used at proof lines 9217-9218. Supplies one of the two inputs of the converse criterion.
\item The converse criterion: quasi-compact plus existence part implies universally closed. Lemma lemma-quasi-compact-existence-universally-closed, used at proof lines 9232-9234. Makes the diagonal universally closed once the criterion is verified for it.
\item A universally closed, locally finite type, separated representable monomorphism is a finite morphism. More on Morphisms, Lemma more-morphisms-lemma-characterize-finite, with Spaces, Lemma spaces-lemma-representable-transformations-property-implication, used at proof lines 9235-9241. Upgrades universal closedness to finiteness of the diagonal.
\item A finite monomorphism is a closed immersion. Morphisms, Lemma morphisms-lemma-finite-monomorphism-closed, used at proof lines 9242-9244. Concludes that the diagonal is a closed immersion, i.e. f is separated.
\end{itemize}\medskip
\ifdefined\hypertarget\hypertarget{sourcepointer-1}{}\fi
\noindent\textit{Source reference.} Original lemma 03KA: focused-seeds/AGP/sources/stacks-snapshot-20261011/tags.txt, source lines 4670--4670. Navigation pointer into the Stacks Project; the standard fact it supplies is applied inside the counted proof span and the argument resting on it stays there.\par\medskip
\ifdefined\hypertarget\hypertarget{sourcepointer-2}{}\fi
\noindent\textit{Source reference.} Original lemma 02LS: focused-seeds/AGP/sources/stacks-snapshot-20261011/tags.txt, source lines 3419--3419. Navigation pointer into the Stacks Project; the standard fact it supplies is applied inside the counted proof span and the argument resting on it stays there.\par\medskip
\ifdefined\hypertarget\hypertarget{sourcepointer-3}{}\fi
\noindent\textit{Source reference.} Original lemma 03BB: focused-seeds/AGP/sources/stacks-snapshot-20261011/tags.txt, source lines 4342--4342. Navigation pointer into the Stacks Project; the standard fact it supplies is applied inside the counted proof span and the argument resting on it stays there.\par\medskip
\ifdefined\hypertarget\hypertarget{sourcepointer-4}{}\fi
\noindent\textit{Source reference.} Original lemma 02YO: focused-seeds/AGP/sources/stacks-snapshot-20261011/tags.txt, source lines 3883--3883. Navigation pointer into the Stacks Project; the standard fact it supplies is applied inside the counted proof span and the argument resting on it stays there.\par\medskip
\begin{lemma}[Valuative criterion separatedness]
\label{lemma-valuative-criterion-separatedness}
Let $S$ be a scheme.
Let $f : X \to Y$ be a morphism of algebraic spaces over $S$.
Assume
\begin{enumerate}
\item the morphism $f$ is quasi-separated, and
\item the morphism $f$ satisfies the uniqueness
part of the valuative criterion.
\end{enumerate}
Then $f$ is separated.
\end{lemma}

\begin{proof}
Assumption (1) means $\Delta_{X/Y}$ is quasi-compact.
We claim the morphism
$\Delta_{X/Y} : X \to X \times_Y X$ satisfies the existence
part of the valuative criterion.
Let a solid commutative diagram
$$
\xymatrix{
\Spec(K) \ar[r] \ar[d] & X \ar[d] \\
\Spec(A) \ar[r] \ar@{-->}[ru] & X \times_Y X
}
$$
be given. The lower right arrow corresponds to a
pair of morphisms $a, b : \Spec(A) \to X$ over $Y$.
By assumption (2) we see that $a = b$. Hence using $a$ as the dotted
arrow works. Hence
Lemma \ref{lemma-quasi-compact-existence-universally-closed}
applies, and we see that $\Delta_{X/Y}$ is universally closed.
Since always $\Delta_{X/Y}$ is locally of finite type and
separated, we conclude from
More on Morphisms, Lemma \ref{more-morphisms-lemma-characterize-finite}
that $\Delta_{X/Y}$ is a finite morphism (also, use the
general principle of
Spaces, Lemma
\ref{spaces-lemma-representable-transformations-property-implication}).
At this point $\Delta_{X/Y}$ is a representable, finite monomorphism,
hence a closed immersion by
Morphisms, Lemma \ref{morphisms-lemma-finite-monomorphism-closed}.
\end{proof}

\end{document}
